\documentclass[11pt,a4paper]{article}
\usepackage[T1]{fontenc}
\usepackage[utf8]{inputenc}
\usepackage[margin=2.3cm,top=2.1cm,bottom=2.0cm]{geometry}
\usepackage{lmodern}
\usepackage{microtype}
\usepackage{xcolor}
\usepackage{enumitem}
\usepackage{titlesec}
\usepackage{tcolorbox}
\usepackage[hidelinks]{hyperref}

\definecolor{ink}{HTML}{16201B}
\definecolor{accent}{HTML}{1A6B3C}
\definecolor{muted}{HTML}{5D6B64}
\definecolor{rule}{HTML}{D7DED9}
\definecolor{soft}{HTML}{F4F7F5}
\definecolor{warn}{HTML}{8A5A12}

\color{ink}
\titleformat{\section}{\large\bfseries\color{accent}}{}{0em}{}
\titlespacing*{\section}{0pt}{13pt}{4pt}
\setlist[itemize]{leftmargin=1.1em,itemsep=2.5pt,topsep=3pt}
\setlength{\parskip}{4pt}
\setlength{\parindent}{0pt}
\pagestyle{empty}

\newtcolorbox{keybox}{colback=soft,colframe=rule,boxrule=0.6pt,
  arc=2pt,left=8pt,right=8pt,top=6pt,bottom=6pt}

\newcommand{\stage}[3]{%
  \vspace{7pt}%
  {\bfseries\color{accent}#1\quad#2}\\[1pt]%
  {\small\color{muted}#3}\par\vspace{2pt}}

\begin{document}

{\LARGE\bfseries What ReqShift does, one stage at a time}\\[3pt]
{\color{muted}The six stages a run passes through, and what each one is for.}

\vspace{4pt}
\begin{keybox}
\textbf{The whole idea in one line.} A language model cannot tell you when it is
wrong, and a solver cannot read prose. So the model reads and proposes, and
everything after that exists to check it. Only one of the six stages guesses
anything.
\end{keybox}

\stage{1}{Reading the requirements}{read}
The text is pulled out of the PDF and the things PDFs break are repaired: page
headers that repeat on every page, words split across a line end, sentences
wrapped over several lines. This matters because every finding later has to point
back at a real sentence in your document, and a sentence broken in half cannot be
quoted or checked.

\stage{2}{Checking the document is in scope}{done}
One cheap question before any money is spent: \emph{does this document say how
many of each thing there can be?} That takes about ten seconds, against twenty
minutes to discover the same thing by extracting first. A document that never
states a quantity has nothing for the analysis to combine, and is refused with the
reason given.

\stage{3}{AI drafts a goal model}{model drafted}
\textbf{The only stage where anything is guessed.} The model reads the whole
document and proposes the participants, what each one has to achieve, how those
goals break into sub-goals, who depends on whom, and how many of each there can
be. Everything after this exists to check this step. It is a draft, and it is
labelled as one wherever it appears.

\stage{4}{Checking the model is well formed}{well formed}
No AI here, only structure. Does every goal have an owner, does every participant
have an overall goal, does anything point at something that is not there. A
malformed draft is caught now rather than after nine audits have been paid for.

\stage{5}{Auditing what each number counts}{the long one}
Nine separate passes, each asking one narrow question, each a fresh call with no
memory of the extraction so it is not anchored to the first answer. This is where
most of the time and most of the cost goes.

\begin{itemize}
\item Does that sentence really state this maximum, or does the number count
      something else? ``The equipment shall have six output ports'' bounds ports,
      not clocks.
\item Can there be \textbf{zero} of this, or did the analysis assume so? That one
      field decides whether there is a finding at all.
\item Are the capacity, demand and timing numbers real?
\item Does each dependency actually exist in the text, or is it merely plausible?
\item Are the AND/OR breakdowns and the overall goals supported by a sentence?
\item What did the extraction \textbf{miss}? Every other pass is blind to an
      omission, and an omission turns a real defect into a clean bill of health.
\end{itemize}

An audit is an improvement to the model, never a gate on it. One that cannot
finish costs its own contribution and nothing else, and the model is saved after
each pass so no failure can waste the work already done.

\stage{6}{The solver proves what follows}{z3, or GoalShift}
The audited model goes to the solver, which works out every build the bounds
allow, checks each one, and groups them into distinct behaviours. Nothing is
sampled and nothing is summarised.

If the banner says \textbf{z3} rather than GoalShift, the model stated durations
or capacities, which the SAT encoding would silently ignore rather than use. The
engine follows the model rather than a dropdown, and the result then carries the
caveat that it is not proved by the encoding the thesis describes. That is
honesty, not an error.

\section*{What comes back}

\textbf{First the gap report:} the quantities the document never states, ranked by
whether their absence can actually break something, each with the sentence it was
read from. Every entry reports an absence, which you can check by searching your
own document.

\textbf{Then the findings:} which builds cannot work, and which participant is
present and provably unable to do its job.

\begin{keybox}
\textbf{Read the marks.} A finding that rests on something the document does not
state is labelled \textbf{unconfirmed} and carries the sentence it depends on. The
solver result is correct for the model as drawn. Whether the model is right about
your specification is the part still open, and the report says which is which
rather than counting them together.
\end{keybox}

\vspace{6pt}
{\footnotesize\color{muted}Model drafted by AI, analysed by GoalShift. The
extraction is a draft for review. The analysis is proved of that draft.}

\end{document}
